\section{Flux d'execution}

\subsection{Lancement principal}

Le flux de lancement courant passe par:

\begin{itemize}
  \item le script \path{APEX/tools/sim/apex_sim_up.sh}, qui construit et lance l'environnement;
  \item ou directement par \path{ros2 launch rc_sim_description apex_sim.launch.py}.
\end{itemize}

Le launch \path{apex_sim.launch.py} charge le scenario, prepare les parametres, lance Gazebo, publie l'etat robot, spawne le vehicule et demarre les ponts necessaires.

\subsection{Sequence runtime}

La sequence typique est la suivante:

\begin{enumerate}
  \item selection du scenario et lecture de \path{apex_sim_scenarios.json};
  \item lancement de Gazebo Sim avec le monde choisi;
  \item generation du robot a partir du Xacro;
  \item publication des transforms par \path{robot_state_publisher};
  \item spawn du modele dans Gazebo;
  \item bridge de \path{/clock}, \path{/apex/sim/scan} et \path{/apex/sim/imu};
  \item inclusion de \path{apex_pipeline.launch.py} en mode simulation;
  \item publication de la commande \path{/apex/cmd_vel_track};
  \item conversion vers \path{/apex/sim/pwm/motor_dc} et \path{/apex/sim/pwm/steering_dc};
  \item conversion finale vers les commandes de joints Gazebo par \path{apex_gz_vehicle_bridge.py}.
\end{enumerate}

\subsection{Role de \texttt{ros\_gz\_bridge}}

Gazebo et ROS 2 utilisent des systemes de communication differents. Le pont \path{ros_gz_bridge} convertit les messages Gazebo en messages ROS 2. Dans cette simulation, il est essentiel pour exposer l'horloge de simulation, les scans LiDAR et l'IMU a la pile APEX.
